\documentclass[11pt,a4paper]{article}

\usepackage[utf8]{inputenc}
\usepackage[T1]{fontenc}
\usepackage[spanish]{babel}
\usepackage{geometry}
\usepackage{graphicx}
\usepackage{array}
\usepackage{tabularx}
\usepackage{makecell}
\usepackage{xcolor}
\usepackage{fancyhdr}
\usepackage{titlesec}
\usepackage{helvet}
\usepackage{ragged2e}
\usepackage{enumitem}
\usepackage{setspace}

\renewcommand{\familydefault}{\sfdefault}
\geometry{
  a4paper,
  left=3.0cm,
  right=3.0cm,
  top=2.35cm,
  bottom=2.35cm,
  headheight=1.9cm,
  headsep=0.45cm
}

\setlength{\parindent}{0pt}
\setlength{\parskip}{0.22em}
\setlength{\tabcolsep}{6pt}
\renewcommand{\arraystretch}{1.08}

\pagestyle{fancy}
\fancyhf{}
\fancyhead[L]{\includegraphics[height=1.35cm]{reporte_tecnico_overleaf_logo.png}}
\fancyhead[R]{\vspace{0.1cm}\small FEIRNNR - Carrera de Computación}
\renewcommand{\headrulewidth}{0.5pt}
\renewcommand{\footrulewidth}{0pt}
\makeatletter
\renewcommand{\headrule}{%
  {\color{gray!65}\hrule width\headwidth height\headrulewidth \vskip-\headrulewidth}%
}
\makeatother

\titleformat{\section}
  {\bfseries\fontsize{15}{17}\selectfont}
  {\thesection.}{0.8em}{}
\titlespacing*{\section}{0pt}{0.7ex}{0.2ex}

\newcolumntype{L}[1]{>{\RaggedRight\arraybackslash}p{#1}}
\newcolumntype{Y}{>{\RaggedRight\arraybackslash}X}

\begin{document}
\thispagestyle{fancy}

\vspace*{0.15cm}

\begin{center}
{\bfseries\fontsize{23}{26}\selectfont
Reporte Técnico de Actividades\\
Práctico-Experimentales Nro. 001\par}
\end{center}

\section{Datos de Identificación del Estudiante y la Práctica}

| Campo                                      | Valor                                                                 |
|---|---|
| **Nombre del estudiante(s)**              | Edison Leonardo Coronel Romero                                       |
| **Asignatura**                            | Desarrollo Basado en Plataformas                                     |
| **Ciclo**                                 | 5 A                                                                  |
| **Unidad**                                | 1                                                                    |
| **Resultado de aprendizaje de la unidad** |                                                                      |
| **Práctica Nro.**                         | 001                                                                  |
| **Título de la Práctica**                 | Implementar un servicio REST con Node.js                             |
| **Nombre del Docente**                    | Edison Leonardo Coronel Romero                                       |
| **Fecha**                                 | Viernes 3 de octubre                                                 |
| **Horario**                               | 07h30 -- 10h30                                                       |
| **Lugar**                                 | Laboratorio Computación aplicada<br>Laboratorio Desarrollo de Software<br>Laboratorio de redes y Sistemas Operativos<br>Laboratorio Virtual<br>EVA<br>Aula |
| **Tiempo planificado en el Sílabo**       | 3 horas                                                              |

\section{Objetivo(s) de la Práctica}
Escriba el objetivo de la práctica (copiado de la guía proporcionada por el docente).

\section{Materiales, Reactivos, Equipos y Herramientas}
Liste los materiales, reactivos, equipos y herramientas utilizados en la práctica, confirmando los de la guía o agregando los adicionales.

\section{Procedimiento / Metodología Ejecutada}
Describa brevemente los pasos que siguió durante la ejecución de la práctica.

\section{Resultados}
Incluya evidencias del trabajo realizado (tablas, gráficos, capturas de pantalla, registros de ejecución, modelos, programas, informes, etc.).

\newpage

\section{Preguntas de Control}
Responda a las preguntas del docente.

\section{Conclusiones}
Redacte las principales conclusiones de la práctica, vinculándolas con el objetivo planteado.

\section{Recomendaciones}
Incluya sugerencias para mejorar la práctica o su aplicación en casos reales.

\section{Bibliografía / Referencias}
Liste las fuentes adicionales consultadas siguiendo el formato IEEE.

\section{Anexos}
Recursos de apoyo para cada una de las secciones anteriores.

\end{document}